% !TEX root = ../main.tex
\section{Transfer Pricing -- interne afregningspriser}\label{thr:transfer-pricing}

\paragraph{Definisjon.} \emph{Transfer pricing} er den prisen som settes når en divisjon i et konsern leverer en vare eller tjeneste internt til en annen divisjon. Prisen påvirker både den enkelte divisjonens regnskapsmessige profit (S2), insentivene til divisjonslederne (S3) og hvilke produksjons- og investeringsbeslutninger som faktisk treffes (S1). Den teoretisk korrekte transfer prisen er lik selgerdivisjonens \emph{opportunity cost} (offeromkostning) ved den interne leveransen. [SLIDE]

\subsection*{Forklaring}

Transfer pricing oppstår nødvendigvis når et konsern desentraliserer ansvar i \thref{responsibility-centers}{profit- eller investmentcentre} med interne leveranser. Hver divisjon har sitt eget P\&L; samme transaksjon må derfor prises som inntekt for selger og kostnad for kjøper. Valget av prismetode bestemmer både \emph{fordelingen} av profit mellom divisjonene og \emph{nivået} av intern handel -- og dermed hvor effektivt konsernet utnytter sine egne ressurser. [SLIDE/BOK]

Pensum skiller mellom to informasjonsregimer som gir vidt forskjellige konklusjoner:

\paragraph{(a) Costless information (symmetrisk, full informasjon).} Toppledelsen kjenner alle marginalkostnader, kapasiteter og opportunity cost-funksjoner. Førstebestes-løsningen er å sette transfer prisen lik selgerdivisjonens \emph{marginal opportunity cost}: hvis det er ledig kapasitet er dette marginalkostnaden ($MC$); hvis kapasiteten er bundet, er det marginalkostnaden \emph{pluss} dekningsbidraget på den fortrengte eksterne salgsenheten. Dette gir konsernoptimal mengde, fordi kjøperdivisjonen utvider intern handel akkurat så lenge dens marginale betalingsvilje overstiger den reelle alternativkostnaden. Resultatet er ekvivalent med MR=MC for konsernet samlet. [SLIDE]

\paragraph{(b) Asymmetrisk informasjon.} I praksis sitter divisjonslederne på \emph{spesifik viden} (jf.\ \thref{specific-vs-general-knowledge}{spesifik vs.\ generell viden}) om egne kostnader, kapasitet og marked. Toppledelsen kan ikke verifisere kostnadsdata uten kostnad. Dette skaper to problemer som rammer alle prisingsmetoder: (i) selgerdivisjonen har insentiv til å \emph{overrapportere} kostnader for å presse opp transfer prisen og dermed sin egen profit; (ii) kjøperdivisjonen kan tilsvarende overdrive interne alternativer for å presse prisen ned. Dette er et klassisk \thref{principal-agent}{prinsipal-agent}-problem på tvers av interne grenser, og det er grunnen til at ingen enkelt metode er optimal. Trade-off-en mellom \emph{decision management} (riktig pris for riktig beslutning) og \emph{decision control} (riktig insentiv for lederen) blir tydelig: presise tall er ikke nødvendigvis insentivkompatible. [SLIDE/BOK]

\subsection*{De fem metodene i pensum}

\paragraph{1. Market-based transfer pricing.} Bruker eksterne markedspriser når det finnes et konkurransedyktig marked for varen. Når markedet er fullkomment, sammenfaller markedspris og opportunity cost, og denne metoden gir både konsernoptimal beslutning og rettferdige insentiver. Krever et reelt og likvid eksternt marked. [SLIDE]

\paragraph{2. Marginal cost (MC) pricing.} Transfer pris $=$ selgerens marginalkostnad. Gir konsernoptimale beslutninger \emph{når kapasiteten er ledig} (opportunity cost $=$ MC). Problem: dekker ikke faste kostnader, så selgerdivisjonen taper penger på interne salg og demotiveres -- og lederen får insentiv til å overrapportere MC. [SLIDE]

\paragraph{3. Full cost pricing.} Transfer pris $=$ MC $+$ allokerte faste kostnader. Enkel å beregne, vanlig i praksis. Problemet er at full cost typisk \emph{overstiger} opportunity cost, slik at intern mengde blir for lav (kjøperen vrir seg mot eksterne leverandører selv når det er suboptimalt for konsernet) -- en intern \thref{double-markup}{double markup}-effekt. [SLIDE]

\paragraph{4. Full cost plus markup (cost-plus).} Full cost $+$ et avtalt prispåslag. Gir selger et garantert dekningsbidrag og dermed bedre insentiver. Forsterker likevel double-markup-problemet og gir lederne enda sterkere insentiv til å skjule reell MC. [SLIDE/BOK]

\paragraph{5. Negotiated transfer pricing.} Divisjonene forhandler prisen seg imellom. Kan ende nær opportunity cost \emph{når begge har omtrent like sterk forhandlingsposisjon}. Resultatet avhenger av forhandlingsevner, ikke av økonomi; tidkrevende; risiko for stadige konflikter som må eskaleres til toppledelsen. [SLIDE]

\paragraph{(Bonus) Dual pricing.} Selger krediteres til markedspris (eller MC $+$ markup), mens kjøper debiteres til marginal cost. Konsernregnskapet rettes så opp internt. Løser insentivproblemet på papiret, men er kostbart, gjør resultatregnskapene mindre transparente og åpner for spill. [BOK]

\subsection*{Kjerneforutsetninger}
\begin{itemize}
\item \textbf{Desentralisering:} Konsernet har minst to enheter med egne P\&L (typisk profit- eller investmentcentre).
\item \textbf{Intern handel:} Det finnes et faktisk varestrøms- eller tjenesteflyt mellom enhetene.
\item \textbf{Målbarhet av kostnad og inntekt:} Kostnader skal kunne henføres til den interne transaksjonen.
\item \textbf{Antakelse om informasjonsregime:} Costless information (a) gir én anbefaling (opportunity cost-pris); asymmetrisk informasjon (b) krever andrebestes-løsninger.
\item \textbf{Eksternt marked (kun market-based):} Det må finnes et likvid og konkurransedyktig eksternt marked for samme vare.
\item \textbf{Kapasitetsforutsetning:} MC-pricing forutsetter ledig kapasitet; ved kapasitetsbinding må prisen heves med opportunity cost av fortrengt eksternt salg.
\end{itemize}

\subsection*{Muligheter}
\begin{itemize}
\item Lar konsernet beholde fordelene ved desentralisering (lokal kunnskap, motivasjon, klare resultatansvar) og samtidig styre intern ressursbruk.
\item Gjør intern verdiskaping synlig i regnskapet og legger grunnlag for divisjonell prestasjonsmåling.
\item Eksponerer interne ineffektiviteter (en divisjon som ikke er konkurransedyktig mot markedet blir tydelig).
\item Brukes også i skatte\-optimalisering (multinasjonale konserner allokerer profit til lavskatteland) -- men dette er regulert av OECDs armlengde-prinsipp.
\end{itemize}

\subsection*{Begrensninger og kritikk}
\begin{itemize}
\item \textbf{Ingen metode er førstebeste under asymmetrisk informasjon.} Alle gir trade-off mellom riktig beslutning og riktig insentiv (decision management vs.\ decision control).
\item \textbf{Gaming og manipulasjon:} Lederne har insentiv til å rapportere kostnader strategisk -- både for å øke egen profit og for å påvirke fremtidige forhandlinger (\emph{ratchet}-lignende effekt).
\item \textbf{Interne konflikter:} Transfer pricing-strider er en av de hyppigste kildene til organisatorisk friksjon i divisjonsstrukturer; eskaleres ofte til toppledelsen.
\item \textbf{Double markup:} Når både selger og kjøper legger på et eget margin, blir samlet pris høyere enn samfunnsoptimal -- intern mengde og samlet profit reduseres.
\item \textbf{Markedspriser kan mangle:} For halvfabrikater og spesialiserte tjenester finnes ofte ikke noe likvid eksternt marked.
\item \textbf{Kapasitetsendringer:} Riktig opportunity cost endres med kapasitetsutnyttelsen, så ``riktig'' transfer pris bør i prinsippet variere over tid -- noe få selskaper får til.
\item \textbf{Skatte- og reguleringsrisiko:} Aggressiv intern prising kan utløse skattetvister (jf.\ armlengdeprinsippet, BEPS).
\end{itemize}

\subsection*{Modell / oversikt}

\begin{center}
\renewcommand{\arraystretch}{1.2}
\begin{tabular}{@{}p{3.4cm}p{4.6cm}p{4.6cm}p{2.0cm}@{}}
\toprule
\textbf{Metode} & \textbf{Fordele} & \textbf{Ulemper} & \textbf{Best n\aa r\dots} \\
\midrule
Market-based & Sammenfaller med opportunity cost; rettferdig; lite gaming & Krever likvid eksternt marked; markedspris kan v\ae re volatil & Eksternt marked finnes \\
Marginal cost & Konsernoptimal mengde ved ledig kapasitet & Dekker ikke faste kostnader; selger taper; MC kan overrapporteres & Ledig kapasitet, MC observerbar \\
Full cost & Enkel; dekker faste kostnader & Overstiger opp.\ cost; for lav intern mengde; double markup & Stabil produksjon, ikke marked \\
Full cost + markup & Garantert margin til selger & Forsterker double markup; gaming-insentiv & Selger trenger insentiv \\
Negotiated & Kan n\ae rme seg opp.\ cost; bruker lokal info & Tidkrevende; avhenger av forhandlingsevne; konflikter & Likeverdig forhandlingsstyrke \\
Dual pricing & Selger og kj\o per ser ulike priser, begge motiveres riktig & Kompleks; skjuler tap; spill mot konsernregnskap & Sm\aa\ avgrensede pilot\-prosjekter \\
\bottomrule
\end{tabular}
\end{center}
\textit{Tabell: De fem (+1) transfer pricing-metodene i pensum, med fordele, ulemper og forutsetninger. [SLIDE -- BSZ kap.\ 17]}

% Slide-figur fra BSZ kap. 17 finnes som EMF og må konverteres til PDF/PNG for å \includegraphics-es. Når konvertert, aktiveres:
% \begin{figure}[h]
%   \centering
%   \includegraphics[width=0.7\linewidth]{BSZ_kap_17/by_slide/slide_NNN/<filnavn>.pdf}
%   \caption{Decentralized firm with two profit centers og transfer pricing -- BSZ kap.\ 17. [SLIDE]}
% \end{figure}

\subsection*{Eksempler}
\begin{itemize}
\item \textbf{Equinor (NO):} Letingsdivisjonen overlater olje og gass til Marketing \& Trading, som videreselger på spotmarkedet. Intern transfer pris er knyttet til \emph{markedspris} (Brent / NBP) -- klassisk market-based metode i et likvid marked. [EKS]
\item \textbf{Norsk Hydro (NO):} Aluminium fra smelteverkene leveres til ekstrudering og valseverk internt. Hydro bruker LME-prisen (London Metal Exchange) som transfer pris -- igjen market-based, fordi et globalt marked finnes. [EKS]
\item \textbf{Statkraft (NO):} Internt strømsalg mellom produksjons- og handelsdivisjon prises mot Nord Pool spot, med systemprisen som referanse. [EKS]
\item \textbf{Mærsk:} Containerleie mellom Ocean-divisjonen og Logistics \& Services prises typisk som \emph{negotiated} eller \emph{cost-plus}, fordi det ikke finnes en god eksternpris for konsernspesifikke containerflåter. Dette har historisk skapt friksjon mellom divisjonene. [EKS/CASE]
\item \textbf{Tax-motivert transfer pricing:} Apple og Starbucks har begge vært i offentlige skattetvister i Europa fordi konserninterne prising av immaterielle rettigheter ble ansett som ikke-armlengde. Viser hvordan transfer pricing også er et juridisk og politisk anliggende, ikke kun et internt styringsverktøy. [EKS]
\item \textbf{Mærsk Power-to-X (case):} Grønn hydrogen og grønn metanol fra PtX-anlegget skal leveres til Mærsks shippingdivisjon som drivstoff. Det finnes ikke noe modent eksternt marked for grønn metanol enn\aa, så market-based metoden faller bort. Cost-plus eller negotiated er realistiske alternativer, men begge har klare svakheter. [CASE]
\end{itemize}

\subsection*{Kobling til søyle og andre teorier}
Søyle: \STwo\ (intern prismåling), \SOne\ (beslutningsrett over interne leveranser) og \SThree\ (insentiv til divisjonsledere). Transfer pricing er det klassiske eksempelet på at de tre søylene må designes samlet -- en god prismetode er ubrukelig hvis insentivstrukturen straffer lederen for å bruke den.

Relaterte teorier: \thref{responsibility-centers}{Responsibility centers} (premisset for at transfer pricing oppstår), \thref{controllability-principle}{Controllability-prinsippet} (lederen må evalueres på beslutninger hun faktisk kan ta), \thref{principal-agent}{Principal-agent} og \thref{agency-costs}{Agency costs} (forklarer asymmetrisk informasjon mellom divisjoner og topp), \thref{specific-vs-general-knowledge}{Spesifik vs.\ generell viden} (hvorfor toppledelsen ikke bare kan diktere riktig pris), \thref{decision-management-control}{Decision management vs.\ decision control} (samme tall kan ikke alltid brukes til begge formål), \thref{double-markup}{Double markup} (fellen ved full cost / cost-plus), \thref{boundaries-of-the-firm}{Boundaries of the firm} (når intern handel ikke fungerer, peker det mot vertikal des-integrasjon).

\paragraph{Brukt i spørsmål:} Q02.
